\section{Conclusions}\label{sec:concl}

A formula-first co-design was carried out for a non-ideal asynchronous Buck
converter supplied at \SI{48}{V} and commanded over \SIrange{1}{46}{V}. Its
search space is the set of Bryson excursions of an LQI cost, so every evaluated
point is a Riccati-optimal law with LQR margins and an operating-polytope
certificate, and every point is implemented exactly as a 2-DOF PIDF that
measures the output voltage only.

Three results carry the paper.

\emph{Theory.} The realization is structural, not only nominal: the derivative
filter of the PIDF, with its pole at the ESR zero, reconstructs the capacitor
voltage along every trajectory of the converter, so the PIDF and the full-state
LQI coincide under saturation, anti-windup and discontinuous conduction and for
any inductor-branch or semiconductor parameters. They differ only by an explicit
feedback on the unmodelled output current, by a mismatch of the output network,
and by sampling. Under sampling the equivalence survives when the
filter is discretized as the observer it is: a triangle-hold update keeps the
sampled PIDF within \EqSampFOH{} of the sampled LQI on a \SI{45}{V} step, where
the usual zero-order-hold update departs by \EqSampZOH{}.

\emph{Optimization.} Seven algorithms applied to the four LQI weights converge
to one optimum; they differ in reliability, pIGWO and pAEABC reaching
feasibility in 28 and 25 of 30 seeds against 13 for GA, and the Gaussian-process
refinement matters most where the global stage is weakest.
\FamConclusion

\emph{Validation.} Switched admissibility is not inherited from an averaged
model: an analytical ripple bound and a steady switched screen both admitted
designs that limit-cycle on the reference mission, because the loop holds
coexisting switching limit cycles selected by the preceding transient. A
\SI{34}{ms} switched transition screen inside the search removes them. On the
reference mission, the retained CMO--LQI--PIDF removes the overshoot of the
classical designs and stays within the \SI{23}{A} device rating; its only miss
is a steady error of \SwEss{}\,mV at \SI{1}{V} against a \SI{25}{mV} limit,
which its full-state twin meets. PI, PID, PIDF, LQR, LQG and LQI designed by their standard rules
exceed at least two specifications by factors of 6 to 30.

The contribution is therefore a design procedure an industrial engineer can
apply with standard tools: choose admissible excursions, solve one Riccati
equation, realize the gain as a PIDF with setpoint weight one and a derivative
pole on the ESR zero, discretize the filter as a capacitor-voltage observer, and
let a small, constrained, multi-fidelity search tune four physically meaningful
numbers. Hardware validation, component-tolerance robustness and line
regulation are the next evidence levels.
